%%%%%%%%%%%%%%%%%%%%%%%%%%%%%%%%%%%%%%
% LaTeX poster template
% Created by Nathaniel Johnston
% August 2009
% http://www.nathanieljohnston.com/2009/08/latex-poster-template/
%%%%%%%%%%%%%%%%%%%%%%%%%%%%%%%%%%%%%%

\documentclass[final]{beamer}
\usepackage[scale=1.24]{beamerposter}
\usepackage{graphicx}			% allows us to import images
\usepackage{bbm,amssymb,mathrsfs}
%\usepackage{natbib}
\usepackage{graphicx,psfrag,dcolumn,amsmath,bm,color}
\usepackage{amssymb}
\usepackage{exscale}
\usepackage{textpos}
\usepackage[tight]{subfigure}
%-----------------------------------------------------------
% Define the column width and poster size
% To set effective sepwid, onecolwid and twocolwid values, first choose how many columns you want and how much separation you want between columns
% The separation I chose is 0.024 and I want 4 columns
% Then set onecolwid to be (1-(4+1)*0.024)/4 = 0.22
% Set twocolwid to be 2*onecolwid + sepwid = 0.464
%-----------------------------------------------------------

\newlength{\sepwid}
\newlength{\onecolwid}
\newlength{\twocolwid}
\newlength{\threecolwid}
\setlength{\paperwidth}{48in}
\setlength{\paperheight}{36in}
\setlength{\sepwid}{0.024\paperwidth}
\setlength{\onecolwid}{0.22\paperwidth}
\setlength{\twocolwid}{0.464\paperwidth}
\setlength{\threecolwid}{0.708\paperwidth}
\setlength{\topmargin}{-0.5in}
\usetheme{confposter}
\usepackage{exscale}

%-----------------------------------------------------------
% The next part fixes a problem with figure numbering. Thanks Nishan!
% When including a figure in your poster, be sure that the commands are typed in the following order:
% \begin{figure}
% \includegraphics[...]{...}
% \caption{...}
% \end{figure}
% That is, put the \caption after the \includegraphics
%-----------------------------------------------------------

\usecaptiontemplate{
\small
\structure{\insertcaptionname~\insertcaptionnumber:}
\insertcaption}

%-----------------------------------------------------------
% Define colours (see beamerthemeconfposter.sty to change these colour definitions)
%-----------------------------------------------------------

\setbeamercolor{block title}{fg=ngreen,bg=white}
\setbeamercolor{block body}{fg=black,bg=white}
\setbeamercolor{block alerted title}{fg=white,bg=dblue!70}
\setbeamercolor{block alerted body}{fg=black,bg=dblue!10}

%-----------------------------------------------------------
% Name and authors of poster/paper/research
%-----------------------------------------------------------

\title{\begin{textblock}{1}(0,0.32)
\includegraphics[scale=0.8]{RevisedPlainLogo.jpg}
\end{textblock}Visualization and Causal Inference of the Mexican Drug War}
\author{Valeria Espinosa and Joseph Kelly \vskip0.5ex \small{vespinos@fas.harvard.edu, kelly2@fas.harvard.edu}}
\institute{Statistics Department, Harvard University}%\\[\medskipamount]


%-----------------------------------------------------------
% Start the poster itself
%-----------------------------------------------------------

\begin{document}
\begin{frame}[t]
  \begin{columns}[t]												% the [t] option aligns the column's content at the top
    \begin{column}{\sepwid}\end{column}			% empty spacer column
    \begin{column}{\onecolwid}
 				\begin{block}{Abstract}
		In Mexico, the presidency of Felipe Calderón (2006-2012) has been characterized for the war against organized crime, raising many questions regarding security and violence. %We attempt to visualize and analyze homicide rates at the municipality level, and link this to information obtained about the association of drug cartels to municipalities.
		The main question of interest is if homicide rates increase significantly after a military intervention. Due to the observational nature of the study we explore the feasibility of causal inference for the data obtained. Some of the challenges are the violation of SUTVA (stable unit treatment value assumption), making ignorability feasible and the partial missingness of key variables. In particular, population information is only observed every 5 years and thus the yearly population of municipalities needs to be imputed.

		There are 
		      \end{block}
      \vskip2ex
      \begin{block}{Objective}
 Our goal is to use the potential outcomes framework to answer the question of interest: \textbf{Have the Military Interventions Increased Violence?}
      \end{block}
		\begin{block}{SUTVA}
			SUTVA - stable unit treatment value assumption
			\begin{itemize}
				\item \textbf{No hidden values of treatments} A broad definition of what ``military intervention'' means in this context helps us think of a two level treatment: receiving a military intervention (defined as ... see paper(that have resulted in deaths?)) or not receiving it.
				\item \textbf{No interference between units} The main idea is to group neighboring regions that have received military interventions in such a way that distances make the ``no interference'' assumption  more reasonable. For treated regions that are side to side were also assessed in terms of neighboring  geographic situation such as lack of highways connecting them %\url{http://mx.kalipedia.com/kalipediamedia/geografia/media/200805/11/geomexico/20080511klpgeogmx_4_Ges_SCO.png} or big mountains between them, or closeness to where the crops or smuggling routes are \url{http://utopiaguatemala.files.wordpress.com/2011/11/mexico_routes_in_466.gif} (can we find such data?).
		We proposed 18 units out of which 5 are eliminated for analysis purposes because the interventions received

			\end{itemize}

	 \end{block}

      \vskip2ex

    \end{column}

    \begin{column}{\sepwid}\end{column}			% empty spacer column
    \begin{column}{\threecolwid}					  % create a three-column-wide column and then we will split it up later
      \begin{block}{Altering Column Spans}
        You can make columns that span multiple other columns relatively easily. Lengths are defined in the template that make columns look normal-ish if you want to use a four-column layout like this poster. If you want to use a different number of columns, you will have to modify those lengths accordingly at the top of the poster.tex file.
        
        In particular, near the top of the TeX file you will see lines that look like:
        \begin{semiverbatim}
          \hskip1ex\\setlength\{\\sepwid\}\{0.024\\paperwidth\}
          
          \hskip1ex\\setlength\{\\onecolwid\}\{0.22\\paperwidth\}
          
          \hskip1ex\\setlength\{\\twocolwid\}\{0.464\\paperwidth\}
          
          \hskip1ex\\setlength\{\\threecolwid\}\{0.708\\paperwidth\}
        \end{semiverbatim}
        
        Set ``sepwid'' to be some small length somewhere near 0.025 (this is the space between columns). Then if $n$ is the number of columns you want, you should set
        \begin{align*}
          \text{onecolwid} & = \frac{1}{n}(1-(n+1)\times\text{sepwid}), \\
          \text{twocolwid} & = 2\times\text{onecolwid} + \text{sepwid}, \\
          \text{threecolwid} & = 3\times\text{onecolwid} + 2\times\text{sepwid}.
        \end{align*}
      \end{block}
      \begin{columns}[t,totalwidth=\threecolwid]	% split up that three-column-wide column
        \begin{column}{\onecolwid}
          \setbeamercolor{block title}{fg=red,bg=white}%frame color
          \setbeamercolor{block body}{fg=black,bg=white}%body color
          \begin{block}{Block Colours}
            For the standard blocks there are two colours; one for the title and one for the block body:\\
            \begin{semiverbatim}
              {\color{red}\\setbeamercolor}\{block title\}\newline \{fg=red,bg=white\}
            \end{semiverbatim}
            \begin{semiverbatim}
              {\color{red}\\setbeamercolor}\{block  body\}\newline \{fg=black,bg=white\}
            \end{semiverbatim}
            The \emph{fg} colour sets the text colour and \emph{bg} sets the background colour.
            For the normal blocks it makes no sense to use a background colour other than white. You \emph{can} change it, but it will look weird!
          \end{block}
        \end{column}
\begin{column}{\onecolwid}
      \begin{alertblock}{Others}
        The ``alert'' block environment looks like this. It also has justified text, but it has a border and a light background to make it stand out. You can create one like so:
        \begin{semiverbatim}
          {\color{red}\\begin}\{{\color{blue}alertblock}\}\{Title\}\newline
          \hskip1ex.......\newline
          {\color{red}\\end}\{{\color{blue}alertblock}\}
        \end{semiverbatim}
      \end{alertblock}
        \end{column}
        \begin{column}{\onecolwid}
          \setbeamercolor{block alerted title}{fg=black,bg=norange}	% frame color
          \setbeamercolor{block alerted body}{fg=black,bg=white}		% body color
          \begin{alertblock}{Alert Block Colours}
            You can similarly modify the colours for alert blocks (but try not to overdo it):\\
            \begin{semiverbatim}
              {\color{red}\\setbeamercolor}\{block title\}\newline \{fg=black,bg=norange\}
            \end{semiverbatim}
            \begin{semiverbatim}
              {\color{red}\\setbeamercolor}\{block  body\}\newline \{fg=black,bg=white\}
            \end{semiverbatim}
          \end{alertblock}        
        %\end{column}
        %\begin{column}{\onecolwid}
          \begin{block}{References}
            Some references and a graphic to show you how it's done:
            \bibliographystyle{plain}
		        \small{\begin{thebibliography}{99}
		        \bibitem{KLPL06} D.~W. Kribs, R. Laflamme, D. Poulin, M. Lesosky, Quantum Inf. \& Comp. \textbf{6} (2006), 383-399.
		        \bibitem{zanardi97} P. Zanardi, M. Rasetti, Phys. Rev. Lett. \textbf{79},  3306 (1997).
		   \end{thebibliography}}

		      \end{block}
        \end{column}
      \end{columns}
      \vskip2.5ex
    \end{column}
  \begin{column}{\sepwid}\end{column}			% empty spacer column
 \end{columns}
\end{frame}
\end{document}
